% !PRINTED 289
\weberpage{289}
%<<CONT
of $\mathfrak{A}_0$ and is consequently expressible rationally through
the roots of the transformation equation. The sum $A$ remains
unchanged under the substitutions of the group $\mathfrak{B}$, and hence
also under $\mathfrak{L}_0$, while under the substitutions of $\mathfrak{A}$, and
therefore also of $\mathfrak{L}$, it takes two different values, namely,
when $b$ runs through the sequence of non-residues,
\[
A = \overset{a}{\Sigma}\varrho^a, \qquad B = \overset{b}{\Sigma}\varrho^b.
\]

$A$ is therefore a function belonging to the group $\mathfrak{L}_0$, and
by its adjunction $\mathfrak{L}$ is reduced to $\mathfrak{L}_0$.

The values of the sums $A$, $B$ are, however, known (Bd. I, \S\,179):
\[
A = \frac{-1 \pm \sqrt{(-1)^{\frac{p-1}{2}}\,p}}{2}, \qquad
B = \frac{-1 \mp \sqrt{(-1)^{\frac{p-1}{2}}\,p}}{2}.
\]

The sign of the root depends on the choice of the root $\varrho$
and can be determined, but does not come into consideration here.
We have therefore the theorem:

\gesperrt{The group of the transformation equation is $\mathfrak{L}_0$,
when $\sqrt{(-1)^{\frac{p-1}{2}}\,p}$ is adjoined to the domain of
rationality.}

The group of the invariant equation can also be derived, without
the division of the elliptic functions, in the following simple
way; in doing so one obtains, however, only the monodromy
group, i.\,e. the group in the field of the rational
functions of $j(\omega)$, without regard to the constants to be
adjoined. We restrict ourselves to the case of a prime degree
$p$ of the transformation, and set
\weq{20}{v_c = j\!\left(\frac{c+\omega}{p}\right), \qquad v_\infty = j(p\omega), \qquad u = j(\omega),}
where $c$ is to be taken with respect to the modulus $p$.

If one applies to $\omega$ a linear substitution $S = \begin{pmatrix}\alpha, & \beta \\ \gamma, & \delta\end{pmatrix}$,
then $j(\omega)$ remains unchanged, and $v_c$ passes into $v_{c'}$. To
find $c'$, one has to seek a second linear substitution
$S' = \begin{pmatrix}\alpha', & \beta' \\ \gamma', & \delta'\end{pmatrix}$, so that
\[
\begin{pmatrix}\alpha', & \beta' \\ \gamma', & \delta'\end{pmatrix}
\begin{pmatrix}p, & 0 \\ c', & 1\end{pmatrix}
=
\begin{pmatrix}p, & 0 \\ c, & 1\end{pmatrix}
\begin{pmatrix}\alpha, & \beta \\ \gamma, & \delta\end{pmatrix}
\]
%CONT>>
